\section{Đọc gì trước, làm gì trước}
\label{sec:start}

\subsection{Danh sách việc khởi động (một trang)}

Làm theo thứ tự. Mỗi dòng có chỗ đánh dấu. Cột cuối chỉ chỗ đọc chi tiết.

\begin{center}
\small
\begin{tabularx}{\textwidth}{@{}c c L l@{}}
\toprule
& \textbf{\#} & \textbf{Việc} & \textbf{Xem} \\
\midrule
\cb & 1 & Mở repo public \code{minhdang900/uit-cs315-lung-nodule-cxr-vlm} trên GitHub; đọc bảng ``Trạng thái và việc cần làm'' đầu \code{code/RUNBOOK.md}. Muốn push thì xin quyền collaborator. & §\ref{sec:git} \\
\cb & 2 & Clone repo một lần vào \code{/content/drive/MyDrive} (RUNBOOK §0, bước ``Lần đầu tiên''). Clone chỗ khác thì đổi \code{REPO}. & §\ref{sec:clone} \\
\cb & 3 & Trên máy cá nhân: tạo venv, chạy \code{pip install -e ".[dev]"} rồi \code{pytest -q}. Phải thấy 79 passed (cần có torch; nếu không, vài test bị skip). & §\ref{sec:env-local} \\
\cb & 4 & Xin quyền tải \code{google/medgemma-4b-it} trên Hugging Face (chấp nhận điều khoản Health AI Developer Foundations). Tạo HF token. Lưu vào Colab Secrets tên \code{HF\_TOKEN}. & §\ref{sec:hf} \\
\cb & 5 & Mua/chuẩn bị Colab Pro (GPU L4 hoặc A100). Không dùng T4. & §\ref{sec:colab} \\
\cb & 6 & Đọc điều khoản 3 bộ dữ liệu. E0 \textbf{đã xong}: xin GVHD chia sẻ thư mục dữ liệu trên Drive, tạo lối tắt \code{cs315-data} trong \code{My Drive}. & §\ref{sec:clone}, §\ref{sec:data} \\
\cb & 7 & Chỉ khi không được chia sẻ: tự tải dữ liệu và chạy 3 script \code{prepare\_*}. Mọi người chạy \code{data\_stats.py} để có \code{\$O/e0\_data\_edges.json}. & §\ref{sec:prepare} \\
\cb & 8 & Phiên GPU đầu tiên: chạy đủ 7 mục kiểm tra bắt buộc. Mục nào hỏng thì dừng. & §\ref{sec:prechecks} \\
\cb & 9 & Chạy E0 → E5 theo đúng thứ tự. Ghi ngân sách vào \code{outputs/budget.md}. & §\ref{sec:exp-run} \\
\cb & 10 & Chạy \code{make\_report.py}, điền các ô đỏ trong báo cáo. & §\ref{sec:fill-report} \\
\bottomrule
\end{tabularx}
\end{center}

\begin{figure}[H]
\centering
\begin{tikzpicture}[node distance=0.45cm and 0.45cm, every node/.style={font=\scriptsize}]
\node[box,minimum width=2.1cm] (a) {Quyền truy cập\\GitHub + HF};
\node[box,right=of a,minimum width=2.1cm] (b) {Clone vào\\Drive, test CPU};
\node[box,right=of b,minimum width=2.1cm] (c) {Tải + chuẩn bị\\dữ liệu (E0)};
\node[obox,right=of c,minimum width=2.1cm] (d) {7 kiểm tra\\bắt buộc (GPU)};
\node[gbox,below=0.9cm of d,minimum width=2.1cm] (e) {E1--E4\\huấn luyện, dự đoán};
\node[gbox,left=of e,minimum width=2.1cm] (f) {make\_report\\+ chọn trên val};
\node[gbox,left=of f,minimum width=2.1cm] (g) {E5 kiểm toán\\che vùng};
\node[box,left=of g,minimum width=2.1cm] (h) {Điền báo cáo,\\slide, release};
\draw[arr] (a)--(b); \draw[arr] (b)--(c); \draw[arr] (c)--(d); \draw[arr] (d)--(e);
\draw[arr] (e)--(f); \draw[arr] (f)--(g); \draw[arr] (g)--(h);
\end{tikzpicture}
\caption{Lộ trình tổng quát. Hộp cam là cổng bắt buộc: chưa qua thì không chạy tiếp.}
\label{fig:roadmap}
\end{figure}

\subsection{Thứ tự đọc tài liệu}

\begin{enumerate}
\item \textbf{Sổ tay này}, §\ref{sec:start}--§\ref{sec:arch}. Mục đích: hiểu bài toán và các khái niệm. Khoảng 2--3 giờ.
\item \textbf{Spec} \code{docs/design/2026-11-12-lung-nodule-cxr-vlm-lora-design.md}. Đây là văn bản có hiệu lực cao nhất.
      Đọc kỹ §5 (phương pháp), §7 (thí nghiệm), §8 (đánh giá) và \textbf{§8.5 (ghi chú triển khai)}.
      Khi sổ tay và spec mâu thuẫn, \textbf{spec đúng}.
\item \textbf{RUNBOOK} \code{code/RUNBOOK.md}. Đây là nguồn duy nhất cho lệnh chạy. Sổ tay chép lệnh từ đây và giải thích thêm.
\item \textbf{Báo cáo} \code{report/main.tex} (PDF: \code{report/main.pdf}). Đọc lướt chương 1--3 để biết văn phong học thuật. Chương 4b là nơi điền kết quả.
\item \textbf{Nhật ký quyết định} \code{docs/decisions.md}. Biết vì sao mọi thứ đang như vậy, trước khi định đổi.
\item \textbf{Tổng quan tài liệu} \code{docs/lit-review/lung-cancer-cxr/}. Chỉ đọc khi cần trích dẫn. Bắt đầu từ \code{00-review-report.md}. File PDF bài báo không có trong git (bản quyền).
\end{enumerate}

\begin{warnbox}
\textbf{Ba quy tắc quan trọng nhất.} (1) Không bao giờ chọn cấu hình hay ngưỡng bằng số trên tập test.
(2) Không commit dữ liệu, trọng số mô hình, thư mục \code{outputs/}, token, hay PDF có bản quyền.
(3) Không tự nghĩ ra cờ dòng lệnh: chép từ RUNBOOK, hoặc đọc \code{argparse} trong \code{code/scripts/}.
\end{warnbox}
